\chapter{系统需求分析与总体设计}

\section{系统需求分析}

结合高速公路团雾监测的业务需求与当前项目实现条件，本文将系统需求归纳为以下四个层面。

在功能层面，系统需要对视频监控画面中的车辆进行实时检测，同时判断当前天气是否为无雾、均匀雾或局地团雾，并给出一个能够反映雾强度的连续变量。仅输出单一标签无法满足联动管控与视频解释需求，因此检测结果、天气类别与雾浓度必须协同呈现。

在数据层面，系统需要在不具备大规模真实团雾样本的前提下完成训练，这要求数据构建方式具备可复用、可扩展、可持续的特点。换言之，不能将系统建立在“必须不断人工采集真实团雾视频并逐帧标注”的高成本前提上。

在算法层面，系统需要尽量减少模型数量与部署链路复杂度。若车辆检测、天气识别和雾浓度估计分别由三个独立模型承担，则推理链会更长、维护成本更高，也不利于部署和答辩展示。因此，统一模型是更合理的总体方向。

在工程层面，系统不仅要能训练，还必须能推理、能导出、能可视化。毕业设计不仅关注模型思想是否成立，还关注完整工程链路是否打通。为此，数据审计、深度预计算、训练脚本、推理脚本、混合推理、视频导出、QAT 与 ONNX 导出都应当纳入总体设计。

\section{总体技术路线}

本文系统总体技术路线可概括为：

\begin{enumerate}[leftmargin=2.5em]
\item 从 UA-DETRAC 中读取晴天交通监控图像与 XML 标注；
\item 使用 MiDaS\_small 对全部样本离线估计深度，并保存深度缓存；
\item 在训练阶段利用深度图和大气散射模型执行在线造雾；
\item 将雾化图像输入统一多任务模型，联合优化检测、雾分类和 beta 回归；
\item 在视频推理阶段，对输入帧执行联合前向，输出天气类别、beta 和车辆检测框；
\item 对 beta 进行时间平滑，并基于其调整显示阈值；
\item 将结果视频输出为可视化演示文件，并保留导出到 ONNX、QAT/INT8 的接口。
\end{enumerate}

这一技术路线兼顾了物理启发、深度学习表示能力和工程部署要求，避免了完全依赖手工特征或完全依赖难获取真实雾数据的极端方案。

\begin{figure}[htbp]
  \centering
  \includegraphics[width=0.95\textwidth]{fig_01_system_overview.png}
  \caption{系统总体框架图}
  \label{fig:system_overview}
\end{figure}

\section{系统模块划分}

从代码实现角度看，本文系统可划分为以下几个核心模块。

\begin{enumerate}[leftmargin=2.5em]
\item 配置模块。位于 \code{src/config.py}，统一管理数据路径、模型参数、训练超参数、物理参数、推理阈值和导出目录，是训练、推理和导出的统一真实配置源。
\item 数据模块。位于 \code{src/data/}。其中 \code{dataset.py} 负责构建多任务训练数据集，\code{depth_estimator.py} 负责深度估计与缓存生成，\code{preparer.py} 则提供离线 YOLO 数据集整理功能。
\item 模型模块。位于 \code{src/model/}。其中 \code{fog_augmentation.py} 负责在线造雾增强，\code{unified_model.py} 负责统一多任务模型结构定义。
\item 训练模块。位于 \code{src/train.py}，负责数据装载、损失计算、联合训练、续训、QAT 和权重保存。
\item 推理模块。位于 \code{src/inference.py}，负责加载统一模型并对视频帧进行推理与可视化。
\item 辅助脚本模块。位于 \code{scripts/}，用于实现数据审计、深度缓存补算、训练前自检、逐行文档生成以及本文后续新增的混合推理脚本等。
\end{enumerate}

\section{系统特点}

本文系统具有以下几个鲜明特点。

其一，真实配置源统一。与许多“文档配置和代码配置并存、易造成版本漂移”的项目不同，本文明确以 \texttt{src/config.py} 作为真实配置入口，使训练、推理和导出均围绕统一参数展开。

其二，训练与推理几何空间保持一致。无论是数据集侧的 letterbox 处理，还是推理时的视频帧预处理，本文均采用一致的几何对齐方式，从而降低训练-推理几何失配带来的性能波动。

其三，系统从一开始就考虑了部署和容错。例如，训练脚本支持缺失深度缓存自动补算；推理系统支持权重自动回退；导出脚本支持 ONNX 输出；量化路径保留了 QAT 与 INT8 转换能力。这使得本项目不仅是一个研究方法，也是一套具备工程闭环特征的原型系统。

其四，系统具有较强的扩展性。后续若要引入真实雾视频监督、加入可见度等级估计、融合更多交通气象传感器数据，均可在当前框架上继续演进，而无需完全推翻已有实现。
